\section{Limitation}
\textcolor{black}{While R-Zero demonstrates strong improvements in reasoning performance across multiple domains, several limitations remain. First, the core mechanism of R-Zero relies on domains where correctness can be objectively verified. The Challenger–Solver co-evolution process depends on clear, deterministic evaluation signals to produce reliable training feedback. Consequently, applying R-Zero to open-ended or subjective tasks, such as creative writing, dialogue, or preference-driven generation, remains difficult, as these tasks lack unambiguous correctness criteria. In addition, R-Zero currently employs specific labeling and verification strategies that may not generalize to all task types. Developing more robust and broadly applicable labeling mechanisms would further expand the range of domains where R-Zero can be effectively applied.}

\section{Conclusion and Future Work}
\label{sec:conclusion}
In this paper, we introduced \method{}, a fully autonomous self-evolving framework that overcomes data dependency by having a Challenger and Solver co-evolve to create a self-generating curriculum.
Our experiments demonstrate that \method{} significantly improves LLM's reasoning capability in multiple domains. Future work could further focus on improving efficiency, exploring more robust labeling techniques, and expanding \method{} to new domains.
% It is crucial to note, however, that the core mechanism of \method{} is currently suited for domains where correctness can be objectively determined.
Extending this self-evolutionary paradigm to open-ended generative tasks, such as creative writing or dialogue, where evaluation is subjective, remains a significant hurdle for future research. 
% We believe \method{} is a significant step towards creating truly self-evolving LLMs.

\section*{Reproducibility Statement}
To ensure the reproducibility of our research, we provide detailed information regarding our data and experimental setup. All datasets used in this work are publicly available; we provide details on data sources, any postprocess steps in Appendix~\ref{app:traininghyper}. A comprehensive list of all hyperparameters for each experiment can be found in a table in Appendix~\ref{app:traininghyper} and Sec.~\ref{sec:experiments}. 